\begin{figure}[t]
    \centering
    \includegraphics[width=\textwidth]{figures/scheme.pdf}
    \vspace{-1.5em}
    \caption{\small
    \textbf{Illustration of \method \vs existing methods.} The language model, pre-trained on texts of length $L$, predicts the $T$th token ($T\gg L$).
    (a) Dense Attention has \(O(T^2)\) time complexity and an increasing cache size. Its performance decreases when the text length exceeds the pre-training text length.
    (b) Window Attention caches the most recent \(L\) tokens' KV. While efficient in inference, performance declines sharply once the starting tokens' keys and values are evicted. (c) Sliding Window with Re-computation rebuilds the KV states from the \(L\) recent tokens for each new token. While it performs well on long texts, its \(O(TL^2)\) complexity, stemming from quadratic attention in context re-computation, makes it considerably slow. (d) \method keeps the \emph{attention sink} (several initial tokens) for stable attention computation, combined with the recent tokens. It's efficient and offers stable performance on extended texts. Perplexities are measured using the Llama-2-13B model on the first book (65K tokens) in the PG-19 test set.}\label{fig:scheme}
    \vspace{-1.5em}
\end{figure}
